\chapter{The Local}
% OUTLINE Ch1 "The Local (topology as 'valid here')". Laws: CITATION (every
% claim traces to name+date; no originals), LENS (device = illustration only;
% this chapter carries exhibit 1 of 2 + the single apparatus citation),
% Vol-1 laws (no Lean identifiers; fence-in-flight; step-not-move).
% Citations landed in-prose: Hausdorff 1914, Kuratowski 1922, Kelley 1955,
% Alexandroff-Hopf 1935 (T0 tradition), Frechet 1906 (limits before metric
% context, one clause). Gate: Kodo per-claim citation test.

Every statement a physicist makes is local. The reading on the bench is valid
near the bench; the field is linear in a small enough region; the temperature
is uniform across a small enough cell; the law holds in the neighborhood where
it was tested and is extended beyond it only on further evidence. Nobody
teaches this discipline as a separate subject, because it is not a doctrine ---
it is the shape of honest reporting. A measurement is an assertion about a
vicinity, and the physicist already reasons in vicinities all day long. This
chapter does one thing: it shows that the mathematics of vicinities was written
down, once and for all, and that nothing in it asks the reader to believe
anything they do not already practice.

The writing-down has a name and a date. In 1914, in the \emph{Grundz\"uge der
Mengenlehre}, Felix Hausdorff codified the idea of a space in terms of
neighborhoods: to every point, a family of sets called its neighborhoods,
subject to a short list of axioms --- every neighborhood of a point contains
the point; the intersection of two neighborhoods contains a neighborhood; any
set containing a neighborhood is one; and inside every neighborhood of a point
sits a neighborhood in which it is a neighborhood of each of its points. That
last clause is the load-bearing one, and the reader should hear it in lab
language: a region of validity contains, around each of its points, a smaller
region of validity. Validity nests. That is the whole axiom, and the physicist
has never worked a day without it.

The modern presentation turns the same content around and speaks of open sets
--- call a set open when it is a neighborhood of each of its points --- and
then the axioms compress to three: the union of any family of open sets is
open, the intersection of finitely many is open, and the whole space and the
empty set are open. The compression is bookkeeping, not new content; the
standard reference for the assembled subject is Kelley's \emph{General
Topology} of 1955, and an equivalent axiomatization through the closure
operation was given by Kuratowski in 1922. One vocabulary note earns its keep
here: a \emph{topology} on a set is nothing but the chosen family of open sets
--- the declared inventory of which regions count as regions of validity ---
and a \emph{topological space} is a set wearing such an inventory. No distance
has been mentioned. That omission is not a gap; it is the subject.

Two workhorse notions now come for free, and both are exactly what their lab
readings suggest.

Continuity first. A map between topological spaces is continuous when the
preimage of every open set is open --- when every demand for output-validity
can be met by a region of input-validity. Read physically: a process is
continuous when, to guarantee the result within a stated tolerance region, it
suffices to control the input within some region. No rate, no bound, no
modulus --- those belong to metric refinements the physicist may add later.
The bare notion is preservation of vicinity, and Hausdorff's framework states
it without a single number.

Limits second. A sequence converges to a point when every neighborhood of that
point, however small, contains all but finitely many terms of the sequence ---
for every region of validity around the target, the sequence eventually enters
and stays. (The idea of framing limits abstractly predates even the
neighborhood axioms: Fr\'echet's 1906 thesis built convergence into abstract
spaces by other means, and Hausdorff's neighborhoods absorbed the lesson.)
The reader should pause on what is absent: no norm measures how far a term is
from the target, and none is needed. Approach is a topological fact ---
membership in ever-smaller declared regions --- before it is ever a numerical
one. Every ball-and-epsilon limit the physicist has computed was secretly this
statement, with the balls supplied by a metric; the statement itself never
needed the ruler, only the regions.

There is one more notion the vicinity language supports, and it is the one a
measuring instrument forces: separation. Given two points of a space, does
some open set contain one and not the other? If yes for every pair --- every
two distinct points are told apart by some region --- the space satisfies what
the trade calls the T0 condition; stronger separation grades follow (each
point closed, T1; disjoint neighborhoods around distinct points, T2, the
condition Hausdorff himself built into his axioms), and the taxonomy is laid
out in the classical literature descending from Alexandroff and Hopf's 1935
\emph{Topologie}. The physical reading is immediate: the open sets are the
space's stock of distinguishing observations, and separation asks whether the
observations can tell the points apart at all.

And when they cannot, mathematics does not avert its eyes; it quotients. Points
that no open set separates are declared one point; the space collapses along
indistinguishability, and the result --- the same observations, no phantom
distinctions --- satisfies T0 by construction. The construction is standard,
it descends from the same classical literature, and it carries Kolmogorov's
name: the T0 identification, in which a space is replaced by the space of
what its own topology can tell apart. Nothing original is being said here;
this is textbook material, and the only reason to slow down over it is that
the reader is about to see the same shape in an unexpected place.

\begin{fence}
An illustration, not an identity. The instrument this series measures --- the
apparatus documented in Volume 1 --- performs, at one load-bearing joint, an
identification of exactly this shape: states that its reading cannot tell
apart are declared one state, and the flagship number is read off the
identified space. That is the T0 identification's \emph{shape}, exhibited in a
working device; it is offered here as an illustration of the construction the
reader has just met, and as nothing more. The device does not thereby do
topology, and this book does not claim it does. Illustration, not identity.
\end{fence}

Take stock of what is now in hand, because the inventory is the chapter's
point. Nearness, without distance. Continuity, without a rate. Limits, without
a norm. Distinguishability, and the discipline of collapsing what cannot be
distinguished. Every one of these is cited, classical, and older than the
reader's textbooks; every one of them formalizes a habit the working physicist
already has; and not one of them required a ruler. The ruler --- the metric,
the inner product, the entire apparatus of lengths --- is famously going to be
the \emph{last} thing this construction adds, chapters from now, and the
reader should begin suspecting already how much geometry can stand without it.

What the vicinity language cannot yet do is hold coordinates. The physicist's
next reflex after ``valid here'' is to lay down axes and measure off
components, and that reflex --- the chart, the observer's coordinates, and the
discipline of changing between observers --- is the next chapter's subject.
This chapter ends where it began, with the reader's own practice: they have
always worked locally; now the word has axioms, a name, and a date.
